\origpage{37}换言之，一个「迭代 $\varphi$-格」$M$ 是一个 $\varphi$-格，使得由 $u$ 诱导的完整同态 $^\tau M\Rightarrow M$ 不呈现 $\tau$-幂零性\footnote{译注：原文作「ne présente \emph{par} de $\tau$-nilpotence」，疑为「\emph{pas}」之误。}，即使模 $\pi_A$ 约化之后也是如此。

\begin{proposition*}
将 $A$ 换为有限全分歧扩张后，$\varphi$-空间 $V$ 具有迭代 $\varphi$-格。
\end{proposition*}

\begin{lemma*}[Drinfeld]
将 $A$ 换为有限全分歧扩张后，$\varphi$-空间 $V$ 具有满足下述条件的格 $M$：
\begin{itemize}
\item $u$ 把 $^\tau M$ 映入 $M$，
\item $\ol u:V^M\to V^M$ 不是幂零的。
\end{itemize}
\end{lemma*}

\begin{proof}
设 $\varphi:V\to V$ 为与 $u:{}^\tau V\overset{\sim}{\longrightarrow}V$ 相伴的 $\tau$-线性态射。

设 $x:V-\{0\}\to\RR$ 为任意一个赋值。

于是
\[ x_\varphi : v\mapsto x_\varphi(v)=\lim_{n\to\infty}\frac{x(\varphi^n(v))}{q^n} \]
也是赋值，并且满足
\[ x_\varphi(\varphi(v))=q\cdot x_\varphi(v),\quad\forall v\neq 0. \]

可以验证，$x_\varphi$ 的取值必然落在某个 $\frac{1}{N}\ZZ$ 之中。

将 $A$ 换为有限全分歧扩张后，可以假设 $x_\varphi$ 取到值 0。

此时 $\{v\in V,\ x_\varphi(v)\ge 0\}=M$ 即为所求。
\end{proof}

\begin{corollary*}
若 $\wh V$ 表示 $V$ 的完备化，则它有非零子空间 $\wh V'\neq 0$ 具有 $\varphi$-格。
\end{corollary*}

\begin{proof}
取引理中的 $M$ 及其完备化 $\wh M$，令
\[ \Big( \bigcap_{n\ge 0}\varphi^n(\wh M) \Big)\otimes_{\wh A_X}\wh K_X=\wh V'. \]
\end{proof}

\begin{proof}[命题的证明]
在 $V$ 中寻找格 $M$，等价于在 $\wh V$ 中寻找格 $\wh M$。

由推论，存在 $\wh V$ 的递增滤过 $(\wh V_s)$，其中 $\wh V_s$ 是维数为 $s\in\ul r^{+}$ 的子空间，使得每个 $\wh V_s/\wh V_{s^-}$ 都具有一个 \origpage{38}$\varphi$-格 $\wh M_s$。设 $\wh M'_s$ 是 $\wh V_s$ 的一个格，它在 $\wh V_s/\wh V_{s^-}$ 中的像是 $\wh M_s$。那么，只要诸 $d_t$（$t\in\ul r^{-}\cap\ul r^{+}$）足够大，
\[ \wh M=\sum_{s\in\ul r^{+}}\Big( \prod_{\substack{t\in\ul r^{+}\\ t<s}}\pi_A^{d_t} \Big)\wh M'_s \]
即为所求。
\end{proof}

设 $M$ 是 $V$ 的一个迭代 $\varphi$-格。

其型是唯一的划分 $\ul r$，使得
\[ d_s\ge 1\iff s\in\ul r^{-}\cap\ul r^{+},\quad 1\le s<r. \]

完整同态
\[ {}^\tau V^M\to V^M\qquad(\text{其中 } V^M=M/\pi_A M) \]
由以下各项构成：
\begin{itemize}
\item $V^M$ 的递增滤过 $(V_s^M)$，满足
\[ \rg V_s^M=s,\quad s\in\ul r^{-}\cap\ul r^{+}, \]
\item ${}^\tau V^M$ 的递减滤过 $(\ol V_s^M)$，满足
\[ {}^\tau V^M={}^\tau V_s^M\oplus\ol V_s^M, \]
\item 同构
\[ \ol V_{s^-}^M/\ol V_s^M\xrightarrow{\ \sim\ }V_s^M/V_{s^-}^M,\quad s\in\ul r^{+}. \]
\end{itemize}

\begin{definition*}
$V^M$ 的子空间 $W$ 称为「好」的，如果
\begin{itemize}
\item 存在 $s\in\ul r^{+}$ 使得 $V_{s^-}^M\subsetneq W\subseteq V_s^M$，
\item 同构 $\ol V_{s^-}^M/\ol V_s^M\xrightarrow{\ \sim\ }V_s^M/V_{s^-}^M$ 诱导出
\[ \ol V_{s^-}^M\cap{}^\tau W\xrightarrow{\ \sim\ }W/V_{s^-}^M. \]
\end{itemize}
\end{definition*}

\origpage{39}\par\medskip\noindent\textbf{迭代 $\varphi$-格的「初等」变换：}\ \emph{设 $M$ 是相对于 $d_1,\ldots,d_{r-1}$ 的迭代 $\varphi$-格，记其型为 $\ul r$。}

\emph{设 $W$ 是 $V^M$ 的一个秩为 $w$ 的「好」子空间。则}
\[ M'=\mathrm{Ker}\,[M\to M/\pi_A M=V^M\to V^M/W] \]
\emph{是相对于 $d'_1,\ldots,d'_{r-1}$ 的迭代 $\varphi$-格，其中当 $s\neq w$ 时 $d'_s=d_s$，而 $d'_w=d_w+1$。$M'$ 的型 $\ul r'$ 由下式给出：}
\[ \ul r'^{+}=\ul r^{+}\cup\{w\}. \]

\begin{theorem*}
任意 $V$ 的迭代 $\varphi$-格 $M$，都可经有限多次初等变换\ul{（朝这一方向或那一方向）}化为一个迭代 $\varphi$-格 $M'$，后者诱导 $\Cht^{r,d,\ol p\le p}(A)$ 上的一点。\footnote{译注：原文作「peut-être transformé」，其中「peut-être」疑为「peut être」之误。}

这个 $M'$ 是唯一的。
\end{theorem*}

回到 $X\otimes\kappa_A$ 上的丛图表
\[ \cE^M\hookrightarrow\cE'^M\hookleftarrow\cE''^M\qquad(\text{一般纤维为 } V^M), \]
它由一个迭代 $\varphi$-格 $M$ 诱导。

有由极大丛诱导的滤过
\[ (\cE_s^M),\quad(\cE_s'^M),\quad(\cE_s''^M)\quad\text{与}\quad(\ol\cE_s^M), \]
以及一般纤维上良定义的同构
\[ (\ol\cE_{s^-}^M/\ol\cE_s^M)\cong(\cE_s''^M/\cE_{s^-}''^M). \]

\begin{lemma*}
对每个 $s\in\ul r^{-}\cup\ul r^{+}$，都有一个处处良定义的态射图表：
\[
\begin{tikzcd}[row sep=small]
 & & \det(\cE_s^M) & \\
\det({}^\tau\cE_s^M) \arrow[r, hook] & \det({}^\tau\cE^M/\ol{\cE}_s^M) \arrow[r, hook] & \det(\cE_s''^M) \arrow[r, hook] & \det(\cE_s'^M) \arrow[hook, from=2-3, to=1-3]
\end{tikzcd}
\]
\end{lemma*}

\origpage{40}此外，$\cE_s'^M/\cE_s^M$ 嵌入 $\cE'^M/\cE^M$，后者维数为 1，并以 $\oo(\wt\cE^M)=\oo(\kappa_A)$ 为支集。

\begin{corollary*}
对每个 $s$，存在 $X$ 上唯一的点 $x(\cE_s^M)$，使得
\[ \det({}^\tau\cE_s^M)\cong\det(\cE_s^M)\,(\oo(\wt\cE^M)-x(\cE_s^M)). \]
这就是秩 $s$ 处的「退化子」。
\end{corollary*}

若 $W$ 是 $V^M$ 的「好」子空间，记 $\cE_w^M$ 为它在 $\cE^M$ 中生成的极大子丛。

\begin{lemma*}
对每个秩为 $w$ 的真「好」子空间 $W\subsetneq V^M$，若 $M'$ 表示 $M$ 经 $W$ 的初等变换，则有两条典范正合列：
\[ 0\to\ol\cE_w^{M'}\to{}^\tau\cE^{M'}\to{}^\tau\cE_w^M\to 0 \]
\[ 0\to{}^\tau\cE_w^M\to{}^\tau\cE^M\to\ol\cE_w^{M'}\to 0 \]
此外，有典范嵌入
\[ \cE_w^{M'}\hookrightarrow\cE_w^M, \]
而商 $\cE_w^M/\cE_w^{M'}$ 在 $\kappa_A$ 上的维数为 0 或 1。
\end{lemma*}

\begin{proof}
\[ M'\subset M\ \text{诱导}\ \cE^{M'}\to\cE^M \]
且
\[ \pi_A M\subset M'\ \text{诱导}\ \cE^M\to\cE^{M'}. \]
甚至有 $\cE(M')=\mathrm{Ker}\,[\cE(M)\to\cE(M)/\cE_w^M]$。注意到，当然有 $\deg\cE^M=\deg\cE^{M'}$；并且由引理，
\[ \deg\cE_s^M-\deg\cE_s^{M'}=0\ \text{或}\ 1,\quad\forall s\in\ul r^{+}\cup\{w\}=\ul r'^{+}. \]
\end{proof}

\begin{proposition*}
(i) 若 $W$ 是 $M$ 的「好」子空间，则有典范同构
\[ \det({}^\tau\cE_w^M)=\det(\cE_w^M)\,(\oo(\wt\cE^M)-x(\cE_w^M)) \]
\origpage{41}其中 $x(\cE_w^M)$ 是唯一的点，满足
\[ x(\cE_w^M)=x(\cE_{s^-}^M)\quad\text{或}\quad x(\cE_s^M)\quad\text{若}\quad V_{s^-}^M\subsetneq W\subsetneq V_s^M. \]

(ii) 对每个 $s\in\ul r^{+}\cup\{w\}$，有
\[ x(\cE_s^{M'})=x(\cE_s^M)\quad\text{若}\quad\deg\cE_s^M=\deg\cE_s^{M'}, \]
\[ x(\cE_s^{M'})=\Frob(x(\cE_s^M))\quad\text{若}\quad\deg\cE_s^M=\deg\cE_s^{M'}+1. \]
\end{proposition*}

\begin{proof}[（ii）的证明概要]
记 $\oo=\oo(\wt\cE^M)=\oo(\wt\cE^{M'})$，$x_s=x(\cE_s^M)$，$x'_s=x(\cE_s^{M'})$。若 $\deg\cE_s^M=\deg\cE_s^{M'}+1$，则存在点 $y$ 与典范同构
\[ \det(\cE_s^M)\cong\det(\cE_s^{M'})(y). \]
与
\[ {}^\tau\det(\cE_s^M)\cong\det(\cE_s^M)(\oo-x_s), \]
\[ {}^\tau\det(\cE_s^{M'})\cong\det(\cE_s^{M'})(\oo-x'_s) \]
相比较，推出
\[ y=x_s\quad\text{且}\quad x'_s=\tau(x_s), \]
或者
\[ y=\tau(y)\quad\text{且}\quad x'_s=x_s. \]
第二种情形可以排除。
\end{proof}

由此得到：

\begin{theorem*}[全局稳定性定理]
设 $\wt\cE,\wt\cF$ 是 $\Cht^{r,\ol p\le p}(K_A)$ 的两个点，\ul{它们在 Hecke 意义下相互对应，}\allowbreak\ul{并延拓为 $\Cht^{r,\ol p\le p}(A)$ 的点}。

那么，$\wt\cE$ 与 $\wt\cF$ 的特殊化的极点、零点与退化子所构成的集合，彼此只相差 Frob 的幂。
\end{theorem*}

\begin{proof}
Hecke 对应只在有限多个位处作用，因此 $\wt\cE$ 与 $\wt\cF$ 有相同的一般纤维 $V$。

\origpage{42}它们的特殊化形如 $\wt\cE^M,\wt\cF^{M'}$，其中 $M,M'$ 是 $V$ 中的两个迭代 $\varphi$-格\footnote{译注：原文作「$M,M'=2$ $\varphi$-réseaux itérés」，其中「$=2$」疑为「deux」之误。}。

由命题（以及 $M,M'$ 可由一列初等变换相联系这一点），只需证明退化子 $x_s=x(\cE_s^M)$ 与 $x'_s=x(\cF_s^M)$\footnote{译注：原文此处夹注「on a donc remplacé $M'$ par $M$」（即把 $M'$ 换成了 $M$）。} 只相差 Frob 的一个幂。

按定义有
\[ {}^\tau\det(\cE_s^M)\cong\det(\cE_s^M)(\oo-x_s), \]
\[ {}^\tau\det(\cF_s^M)\cong\det(\cF_s^M)(\oo-x'_s), \]
又有
\[ \det(\cF_s^M)\cong\det(\cE_s^M)\Big( \sum_\iota m_\iota x_\iota \Big), \]
从而
\[ (x'_s-x_s)+\sum_\iota m_\iota(\tau(x_\iota)-x_\iota)=0, \]
以及
\[ x'_s\in\tau^{\ZZ}(x_s). \]
\end{proof}
